%% AXIOM v1.0 — Autonomous Value Protocol
%% White Paper — Draft

\documentclass[11pt,a4paper]{article}

\usepackage[utf8]{inputenc}
\usepackage[T1]{fontenc}
\usepackage{amsmath,amssymb,amsthm}
\usepackage{mathtools}
\usepackage{geometry}
\usepackage{hyperref}
\usepackage{booktabs}
\usepackage{array}
\usepackage{enumerate}
\usepackage{xcolor}
\usepackage{pgfplots}
\pgfplotsset{compat=1.18}

\geometry{margin=1in}

\theoremstyle{plain}
\newtheorem{theorem}{Theorem}[section]
\newtheorem{lemma}[theorem]{Lemma}
\newtheorem{proposition}[theorem]{Proposition}
\newtheorem{corollary}[theorem]{Corollary}

\theoremstyle{definition}
\newtheorem{definition}[theorem]{Definition}
\newtheorem{example}[theorem]{Example}

\theoremstyle{remark}
\newtheorem{remark}[theorem]{Remark}

\newcommand{\PVC}{\mathrm{PVC}}
\newcommand{\PVCS}{\mathrm{PVCS}}
\newcommand{\NVC}{\mathrm{NVC}}
\newcommand{\PVE}{\mathrm{PVE}}
\newcommand{\AS}{\mathrm{AS}}
\newcommand{\Rid}{R_{\mathrm{id}}}
\newcommand{\Rstrat}{R_{\mathrm{strat}}}
\newcommand{\E}{\mathbb{E}}
\newcommand{\R}{\mathbb{R}}
\newcommand{\N}{\mathbb{N}}

\title{%
  \textbf{AXIOM v1.0}\\[0.5em]
  \large A Formal Protocol for Capital Allocation\\
  via Proof of Value Creation
}

\author{%
  Autonomous Value Protocol Specification\\[0.3em]
  \small \texttt{https://github.com/haynbroit-alt/cpo}
}

\date{2026 — Draft v0.1}

\begin{document}

\maketitle

\begin{abstract}
We present AXIOM, a decentralised economic protocol for allocating capital to
autonomous agents based on verifiable value creation rather than revenue or
speculation. The central primitive is the \emph{Proof of Value Creation} (PVC),
a formally defined, cryptographically attested record of net value produced by
an agent within a bounded, auditable execution context. Capital allocation
follows an \emph{AxiomScore} — a monotone, risk-adjusted function of
accumulated PVC, identity reputation, and strategy reputation. We prove that
the slashing mechanism makes fraud economically irrational under mild
assumptions, and that the protocol converges to a stable allocation fixed point
when agent populations satisfy a diversity condition. AXIOM is designed to
compose with execution-layer verification protocols (such as Proof
Protocol / CPO~\cite{cpp2025}) via shared proof formats, without structural
dependency.
\end{abstract}

\tableofcontents
\newpage

% ---------------------------------------------------------------------------
\section{Introduction}
% ---------------------------------------------------------------------------

Autonomous agents — software systems that execute economic strategies without
continuous human oversight — are increasingly capable of generating value
through trading, arbitrage, digital commerce, and computation. Yet capital
allocation to such agents remains primitive: investors rely on audited backtests,
historical returns, or reputational proxies. These signals are expensive to
verify, easy to manipulate, and structurally unable to distinguish genuine value
creation from variance extraction or fraud.

AXIOM addresses a single core question:

\begin{quote}
\emph{How should a decentralised protocol allocate capital to agents whose
outputs are measurable, based solely on verifiable proof that value was
actually created?}
\end{quote}

The answer requires three separable components:

\begin{enumerate}[(1)]
  \item A \textbf{definition} of value creation that is objective, computable,
        and resistant to gaming.
  \item A \textbf{verification mechanism} that confirms claimed value creation
        without requiring trust in any single party.
  \item An \textbf{allocation function} that maps verified value creation to
        capital, with incentive properties that make fraud irrational.
\end{enumerate}

This document provides formal treatments of all three.

\subsection{Relationship with Proof Protocol / CPO}

AXIOM is independent of any specific execution-layer technology but is designed
to compose with Proof Protocol~\cite{cpp2025}. The distinction is precise:

\begin{itemize}
  \item \textbf{Proof Protocol / CPO}: proves that a computation was executed
        correctly, deterministically, in a bounded environment (what happened).
  \item \textbf{AXIOM}: decides how to reward and finance agents based on the
        economic value of those executions (what it was worth).
\end{itemize}

A CPO serves as the \emph{proof} component of a PVC: it attests that the code
producing the reported output genuinely ran in a sandboxed environment and
produced the claimed result. AXIOM treats CPOs as an opaque, verifiable
attestation primitive.

% ---------------------------------------------------------------------------
\section{Formal Model}
% ---------------------------------------------------------------------------

\subsection{Agents}

\begin{definition}[Agent]
An agent $a$ is a tuple
\[
  a = \bigl(\mathrm{id}_a,\; \mathrm{stake}_a,\; \mathcal{H}_a,\;
             \Rid(a),\; \Rstrat(a),\; s_a\bigr)
\]
where:
\begin{itemize}
  \item $\mathrm{id}_a \in \{0,1\}^{256}$: stable public identity (e.g.\ public key hash);
  \item $\mathrm{stake}_a \in \R_{\ge 0}$: capital deposited as collateral;
  \item $\mathcal{H}_a$: ordered history of finalised PVCs attributed to $a$;
  \item $\Rid(a,t) \in [0,1]$: identity reputation at time $t$;
  \item $\Rstrat(a,t) \in [0,1]$: strategy reputation at time $t$;
  \item $s_a \in \mathcal{S}$: current strategy class.
\end{itemize}
\end{definition}

Let $\mathcal{S} = \{\text{trading, arbitrage, optimisation, marketing,
automation, digital-commerce}\}$ be the finite set of recognised strategy
classes.

\subsection{Proof of Value Creation}

\begin{definition}[PVC]
A \emph{Proof of Value Creation} for agent $a$ over interval $[t_0, t_1]$ is a
tuple
\[
  \PVC = \bigl(\mathrm{id}_a,\; s_a,\; V_0,\; V_1,\; C,\; \tau,\; \sigma,\;
               \pi\bigr)
\]
where:
\begin{itemize}
  \item $V_0, V_1 \in \R_{\ge 0}$: observable value at start and end of the
        period, measured in a common numeraire;
  \item $C \in \R_{>0}$: capital allocated by the protocol during $[t_0, t_1]$;
  \item $\tau = t_1 - t_0 > 0$: duration;
  \item $\sigma \in \R_{>0}$: realised volatility of the agent's value process
        over $[t_0, t_1]$;
  \item $\pi$: cryptographic execution proof (e.g.\ a CPO) attesting that the
        reported $(V_0, V_1, \sigma)$ derive from a verifiably executed strategy.
\end{itemize}
\end{definition}

\begin{definition}[Net Value Created]
\label{def:nvc}
Let $r_f \ge 0$ be the risk-free rate. The \emph{net value created} by a PVC is
\[
  \NVC(\PVC) \;=\; (V_1 - V_0) \;-\; C \cdot r_f \cdot \tau.
\]
\end{definition}

This definition charges the agent for the opportunity cost of the capital
allocated. An agent that merely matches the risk-free return produces
$\NVC = 0$.

\begin{definition}[PVC Score]
\label{def:pvcs}
The \emph{PVC Score} (risk-adjusted) is
\[
  \PVCS(\PVC)
  \;=\;
  \frac{\NVC(\PVC)}{C \cdot \sigma \cdot \sqrt{\tau}}.
\]
\end{definition}

The denominator $C \cdot \sigma \cdot \sqrt{\tau}$ is proportional to the
standard deviation of a position of size $C$ with volatility $\sigma$ over
$\sqrt{\tau}$ time units — an Itô-style risk normalisation. $\PVCS$ is
dimensionless and comparable across agents regardless of strategy size or
duration.

\begin{remark}
$\PVCS$ is structurally identical to a Sharpe ratio with the risk-free rate
subtracted, scaled by capital. A positive $\PVCS$ indicates that the agent
created value beyond the opportunity cost \emph{and} after adjusting for risk.
\end{remark}

\subsection{Protocol Value Efficiency}

\begin{definition}[Protocol Value Efficiency]
At time $T$, the \emph{Protocol Value Efficiency} is
\[
  \PVE(T)
  \;=\;
  \frac{\displaystyle\sum_{a \in \mathcal{A}} \sum_{\PVC_i \in \mathcal{H}_a} \NVC(\PVC_i)}
       {\displaystyle\sum_{a \in \mathcal{A}} \int_0^T C_a(t)\,\mathrm{d}t}.
\]
\end{definition}

$\PVE > 0$ is the necessary condition for the protocol to be economically
self-sustaining. The long-term goal of AXIOM governance is to keep $\PVE$
strictly positive.

% ---------------------------------------------------------------------------
\section{Reputation}
% ---------------------------------------------------------------------------

\subsection{Identity Reputation}

Identity reputation $\Rid$ measures the trustworthiness of an agent identity
over time, independently of strategy performance.

\begin{definition}[Identity Reputation Update]
\label{def:rid}
Let $b_t \in [0,1]$ be a behavior score at time $t$ (1 = fully honest,
0 = confirmed fraud). Then
\[
  \Rid(a, t+1) \;=\; (1 - \lambda_{\mathrm{id}}) \cdot \Rid(a, t)
                    \;+\; \lambda_{\mathrm{id}} \cdot b_t
\]
with $\lambda_{\mathrm{id}} \in (0,1)$ a protocol parameter controlling the
update rate.
\end{definition}

Properties of $\Rid$:
\begin{itemize}
  \item \textit{Monotone recovery}: sustained honest behavior drives $\Rid \to 1$.
  \item \textit{Irreversibility of fraud}: a single $b_t = 0$ event pulls $\Rid$
        downward regardless of prior history.
  \item \textit{Sybil penalty}: a freshly created identity starts at $\Rid = 0$
        and must accumulate history before receiving significant allocation.
\end{itemize}

\subsection{Strategy Reputation}

Strategy reputation $\Rstrat$ measures performance within a specific strategy
class.

\begin{definition}[Strategy Reputation Update]
\label{def:rstrat}
Upon finalisation of $\PVC_i$:
\[
  \Rstrat(a, t+1) \;=\; (1 - \lambda_s) \cdot \Rstrat(a, t)
                       \;+\; \lambda_s \cdot \Phi\bigl(\PVCS(\PVC_i)\bigr)
\]
where $\Phi : \R \to [0,1]$ is a sigmoid squashing function and
$\lambda_s \in (0,1)$ is the strategy update rate.
\end{definition}

\subsection{Strategy Transition and Reputation Transfer}

When agent $a$ switches from strategy class $s$ to $s'$, its accumulated
strategy reputation transfers partially according to a \emph{proximity matrix}.

\begin{definition}[Strategy Proximity Matrix]
$M \in [0,1]^{|\mathcal{S}| \times |\mathcal{S}|}$ is a symmetric matrix where
$M[s, s']$ reflects the structural similarity between strategy classes $s$ and
$s'$. Diagonal entries satisfy $M[s,s] = 1$.
\end{definition}

\begin{definition}[Reputation Transfer]
Upon strategy change from $s$ to $s'$:
\[
  \Rstrat'(a) \;=\; \Rstrat(a) \cdot M[s, s'].
\]
\end{definition}

An agent switching between structurally similar strategies (e.g.\ trading
$\to$ arbitrage) retains most reputation. Switching to a structurally distant
class (e.g.\ trading $\to$ digital-commerce) forfeits most accumulated strategy
reputation, preventing reputation laundering.

% ---------------------------------------------------------------------------
\section{Capital Allocation}
% ---------------------------------------------------------------------------

\subsection{AxiomScore}

\begin{definition}[AxiomScore]
\label{def:axiom-score}
The AxiomScore of agent $a$ at time $t$ is
\[
  \AS(a, t)
  \;=\;
  \Biggl[\sum_{i} \PVCS(\PVC_i) \cdot e^{-\kappa(t - t_i)}\Biggr]^{\!\alpha}
  \cdot
  \Rid(a, t)^{\beta}
  \cdot
  \Rstrat(a, t)^{\gamma}
  \cdot
  \bigl(1 + \mathrm{VaR}_{\delta}(a, t)\bigr)^{-\delta}
\]
where:
\begin{itemize}
  \item $e^{-\kappa(t - t_i)}$: exponential decay with half-life $\ln 2 / \kappa$,
        so that recent PVCs carry more weight;
  \item $\alpha, \beta, \gamma, \delta > 0$: protocol parameters;
  \item $\mathrm{VaR}_{\delta}(a,t)$: the $\delta$-quantile Value-at-Risk of
        agent $a$'s strategy, estimated from recent history.
\end{itemize}
\end{definition}

The four factors encode:
\begin{enumerate}[(1)]
  \item \textit{Performance}: recency-weighted sum of risk-adjusted PVC scores;
  \item \textit{Identity trust}: long-term honesty signal;
  \item \textit{Strategy mastery}: competence within the current strategy class;
  \item \textit{Tail-risk penalty}: penalises high-variance strategies that
        impose systemic risk on the pool.
\end{enumerate}

\subsection{Allocation Function}

\begin{definition}[Capital Allocation]
\label{def:allocation}
The capital allocated to agent $a$ at rebalancing epoch $t$ is
\[
  C_a(t)
  \;=\;
  C_{\mathrm{total}}(t)
  \cdot
  \frac{\AS(a, t)^{\eta}}{\displaystyle\sum_{b \in \mathcal{A}} \AS(b, t)^{\eta}}
\]
where $\eta \ge 1$ is a concentration parameter. Higher $\eta$ produces a
winner-takes-more distribution; $\eta = 1$ is proportional.
\end{definition}

\begin{proposition}[Progressive allocation]
For any agent $a$ with $\AS(a, 0) = 0$, the allocation $C_a(t) \to 0$ as
all other agents maintain positive scores. New agents receive near-zero
allocation until they accumulate verified PVC history.
\end{proposition}

\begin{proof}
Direct from Definition~\ref{def:allocation}: if $\AS(a,0) = 0$ then the
numerator is 0 while the denominator is strictly positive.
\end{proof}

\subsection{Systemic Risk Constraint}

To limit correlated drawdown across the pool, the protocol imposes:

\begin{definition}[Correlation Constraint]
Let $\rho_{ab}$ be the pairwise return correlation between agents $a$ and $b$.
The portfolio concentration constraint requires
\[
  \sum_{a,b:\,\rho_{ab} > \rho^*} C_a(t) \cdot C_b(t)
  \;\le\;
  \theta \cdot C_{\mathrm{total}}(t)^2
\]
for parameters $\rho^* \in (0,1)$ and $\theta \in (0,1)$. Agents violating this
constraint have their allocation capped and the surplus redistributed.
\end{definition}

% ---------------------------------------------------------------------------
\section{Validation and Slashing}
% ---------------------------------------------------------------------------

\subsection{PVC Lifecycle}

Each submitted PVC traverses the following state machine:
\[
  \text{Submitted}
  \;\to\;
  \text{Primary Validated}
  \;\to\;
  \text{Challenge Window}
  \;\to\;
  \begin{cases}
    \text{Finalised} & \text{(no valid challenge)} \\
    \text{Audited}   & \text{(challenge raised)}
  \end{cases}
  \;\to\;
  \begin{cases}
    \text{Accepted} \\
    \text{Rejected}
  \end{cases}
\]

Validators stake capital before participation. Any party may raise a challenge
during the Challenge Window by depositing a challenge bond. An audit panel
(sampled randomly from staked auditors) adjudicates disputes.

\subsection{Slashing Rules}

\begin{definition}[Slashing]
\begin{itemize}
  \item \textbf{Agent fraud} (submitting a false PVC): slash $\mathrm{stake}_a$,
        set $b_t = 0$, permanently flag $\mathrm{id}_a$.
  \item \textbf{Validator fault} (approving a false PVC): slash validator stake.
  \item \textbf{Auditor reward}: honest auditors who correctly identify fraud
        receive a share of slashed collateral.
  \item \textbf{Frivolous challenge}: challenger loses bond if the PVC is
        upheld after audit.
\end{itemize}
\end{definition}

\subsection{Economic Security Theorem}

\begin{theorem}[Fraud Irrationality]
\label{thm:fraud}
Let $p \in (0,1)$ be the probability that a fraudulent PVC is detected during
the Challenge Window or audit. Let $G > 0$ be the expected gain from a
fraudulent PVC (capital received before detection). Then fraud is strictly
irrational if and only if
\[
  \mathrm{stake}_a \;>\; \frac{G}{p}.
\]
\end{theorem}

\begin{proof}
The expected utility of attempting fraud for a risk-neutral agent is
\[
  \E[\text{fraud}]
  \;=\;
  (1-p) \cdot G \;-\; p \cdot \mathrm{stake}_a.
\]
Setting $\E[\text{fraud}] < 0$ and solving:
\[
  (1-p) G < p \cdot \mathrm{stake}_a
  \;\iff\;
  G - pG < p \cdot \mathrm{stake}_a
  \;\iff\;
  G < p(\mathrm{stake}_a + G)
\]
which simplifies to $\mathrm{stake}_a > G/p - G = G(1-p)/p$.
For $p > 0.5$ the right-hand side is less than $G$, so a stake matching
expected gain suffices. For $p \to 1$ the required stake approaches 0;
for $p \to 0$ the required stake diverges (the protocol relies on detection).
The protocol sets minimum stake $S_{\min}$ such that
$S_{\min} \ge G_{\max}(1-p)/p$ where $G_{\max}$ is the maximum extractable
gain per PVC epoch, making fraud irrational for all rational agents.
\end{proof}

\begin{corollary}[Validator Collusion Cost]
For $k$ validators to collude and approve a fraudulent PVC, the expected loss
per validator is $p_k \cdot \mathrm{stake}_v$ where $p_k$ is the per-validator
detection probability given $k$-collusion. Since auditors are sampled
independently, $p_k$ decreases slowly with $k$ (sublinear), meaning the
marginal cost of recruiting an additional colluder remains high. Collusion of
$k > \lfloor |\mathcal{V}| / 3 \rfloor$ validators is required to override
audit, giving Byzantine fault tolerance analogous to BFT consensus.
\end{corollary}

% ---------------------------------------------------------------------------
\section{Attack Analysis}
% ---------------------------------------------------------------------------

\subsection{Sybil Attacks}

An adversary creates $n$ fresh identities to dilute the pool or acquire
disproportionate allocation.

\textbf{Defense}: Each fresh identity starts with $\Rid = 0$ and $\Rstrat = 0$,
yielding $\AS = 0$ (from Definition~\ref{def:axiom-score}). Allocation is
therefore zero until PVC history is accumulated. The cost of building
reputation on $n$ identities is $n$ times the cost of building it on one,
making Sybil attacks linearly expensive.

\subsection{Fake PVC Attacks}

An agent fabricates $(V_0, V_1, \sigma)$ without genuine execution.

\textbf{Defense}: The execution proof $\pi$ (e.g.\ a CPO from Proof Protocol)
is a non-forgeable cryptographic attestation that the reported values derive
from a verifiably executed computation. A fake PVC lacking a valid $\pi$ is
rejected at primary validation. A PVC with a forged $\pi$ requires breaking
the underlying signature scheme (Ed25519), which is computationally
infeasible under standard assumptions.

\subsection{Reputation Laundering}

An agent with high $\Rstrat$ in one class switches to a distant class, exploiting
accumulated reputation to receive elevated allocation before establishing
performance.

\textbf{Defense}: Strategy transition multiplies reputation by $M[s, s'] \ll 1$
for distant classes (Definition~\ref{def:rstrat}). Allocation in the new class
is therefore initially low, matching the genuinely low demonstrated competence.

\subsection{Oracle Manipulation}

An agent manipulates $V_0$ or $V_1$ measurements at the data-source level.

\textbf{Defense}: PVC values must be derived from on-chain or third-party-attested
sources enumerated in the protocol registry. The execution proof $\pi$ commits
to the specific data source used. Disputes over oracle values trigger an audit
of the source attestation, not merely the computation.

\subsection{Correlation Attack}

An adversary controls many agents with identical strategies, creating concentrated
systemic risk in the pool while each individual agent appears low-risk.

\textbf{Defense}: The correlation constraint (Section 4.3) explicitly limits
the aggregate allocation to correlated agent clusters. Agents above the
threshold have allocation capped regardless of individual AxiomScore.

% ---------------------------------------------------------------------------
\section{Governance}
% ---------------------------------------------------------------------------

Governance in AXIOM is deliberately limited to parameter adjustments. The
immutability boundary is:

\begin{itemize}
  \item \textbf{Immutable}: PVC history, finalised allocations, agent identity
        records, slashing events.
  \item \textbf{Governable}: $\alpha, \beta, \gamma, \delta, \eta, \kappa$
        (score parameters); $\rho^*, \theta$ (risk constraints); $\lambda_s,
        \lambda_{\mathrm{id}}$ (reputation update rates); $r_f$ (reference
        rate); strategy class definitions; minimum stake requirements.
\end{itemize}

Governance votes are weighted by $\Rid$ of the voter, not raw token stake,
to prevent governance capture by capital-rich new entrants.

% ---------------------------------------------------------------------------
\section{MVP Specification}
% ---------------------------------------------------------------------------

The MVP instantiates AXIOM with the minimal viable parameter set:

\begin{itemize}
  \item \textbf{Agent class}: crypto trading only ($|\mathcal{S}| = 1$).
  \item \textbf{PVC proof}: on-chain trade records (DEX transaction hashes),
        optionally augmented with CPO execution proofs for strategy code.
  \item \textbf{Validation}: $k = 3$ validators per PVC, majority vote.
  \item \textbf{Challenge Window}: 48 hours.
  \item \textbf{Allocation}: single capital pool, daily rebalancing.
  \item \textbf{Reputation}: $\Rid$ only; $\Rstrat \equiv 1$ (single strategy
        class, no switching).
  \item \textbf{Parameters}: $\alpha = 1, \beta = 0.5, \gamma = 0$ (disabled),
        $\delta = 0.5, \eta = 1$ (proportional), $\kappa = \ln 2 / 30$ (30-day
        half-life), $r_f = $ annualised reference rate.
\end{itemize}

The MVP objective is to demonstrate that a capital pool allocated by AxiomScore
achieves $\PVE > 0$ over a 90-day trial.

% ---------------------------------------------------------------------------
\section{Theoretical Properties}
% ---------------------------------------------------------------------------

\begin{theorem}[Allocation Fixed Point]
\label{thm:fixedpoint}
Suppose the agent population $\mathcal{A}$ is fixed and all agents maintain
constant $\PVCS$. Then the allocation sequence $\{C_a(t)\}_{t \ge 0}$ converges
to a unique fixed point $C^*$ determined by the stationary AxiomScores.
\end{theorem}

\begin{proof}[Proof sketch]
At stationarity, $\Rid(a,t)$ and $\Rstrat(a,t)$ converge to fixed points
$\bar{R}_{\mathrm{id}}(a)$ and $\bar{R}_{\mathrm{strat}}(a)$ (geometric
series limits of the exponential moving average updates in
Definitions~\ref{def:rid}--\ref{def:rstrat}). The sum
$\sum_i \PVCS(\PVC_i) e^{-\kappa(t-t_i)}$ also converges to a constant
under stationary $\PVCS$. Therefore $\AS(a)$ converges, and $C_a^* =
C_{\mathrm{total}} \cdot \AS(a)^\eta / \sum_b \AS(b)^\eta$ is a well-defined
fixed point.
\end{proof}

\begin{theorem}[PVE Positivity under Fraud Irrationality]
\label{thm:pve}
If the minimum stake condition of Theorem~\ref{thm:fraud} holds for all agents
and validators, and at least one agent has strictly positive true $\PVCS$,
then $\E[\PVE(T)] > 0$ for all $T > 0$.
\end{theorem}

\begin{proof}[Proof sketch]
Under fraud irrationality, all submitted PVCs that pass validation have
$\NVC > 0$ in expectation (negative-NVC submissions are irrational since
they waste stake without extracting gain). By monotonicity of the allocation
function, positive-NVC agents receive positive allocation, and the expected
pool-level NVC is strictly positive. Since denominator $\sum_a \int C_a(t)\,dt$
is finite, $\E[\PVE(T)] > 0$.
\end{proof}

% ---------------------------------------------------------------------------
\section{Relation to Proof Protocol / CPO}
% ---------------------------------------------------------------------------

Table~\ref{tab:comparison} summarises the complementary roles.

\begin{table}[h]
\centering
\begin{tabular}{lll}
\toprule
& \textbf{Proof Protocol / CPO} & \textbf{AXIOM} \\
\midrule
Question answered & Was this computation correct? & Was this computation valuable? \\
Core primitive    & Computational Proof Object (CPO) & Proof of Value Creation (PVC) \\
Verification      & Deterministic re-execution & Economic audit + slashing \\
Output            & Signed execution record & Capital allocation \\
Scope             & Any bounded computation & Agent economic strategies \\
\bottomrule
\end{tabular}
\caption{Layered separation of concerns.}
\label{tab:comparison}
\end{table}

A PVC $\pi$ field pointing to a CPO allows any party to verify not only
that the agent's reported value was produced honestly (AXIOM audit) but also
that the underlying computation ran in a provably isolated environment
(Proof Protocol). The two layers are independent: AXIOM can function without
CPO attestations (relying on on-chain proofs), and Proof Protocol is useful
without AXIOM (in any verified computation context).

% ---------------------------------------------------------------------------
\section{Discussion and Future Work}
% ---------------------------------------------------------------------------

Several important questions remain open:

\paragraph{Value definition for non-financial agents.}
Definition~\ref{def:nvc} assumes a monetary numeraire. Extending AXIOM to
agents producing non-monetary value (information, coordination, content) requires
a generalised value function $V : \text{outcome} \to \R$ that is
auditable and manipulation-resistant. This is an open research problem
in mechanism design.

\paragraph{Calibration of $\alpha, \beta, \gamma, \delta$.}
The AxiomScore exponents determine the relative importance of performance,
identity, strategy, and risk. Optimal calibration requires empirical data
from the MVP deployment.

\paragraph{Cross-pool allocation.}
When multiple AXIOM pools coexist with different risk profiles, inter-pool
capital flows require a meta-allocation mechanism. We conjecture that a
hierarchical AxiomScore applying pool-level PVE as an additional factor
would suffice, but this requires formal treatment.

\paragraph{Dynamic strategy classes.}
The current model treats $\mathcal{S}$ as fixed. In practice, new strategy
classes emerge as AI capabilities advance. A governance process for introducing
new classes with appropriate proximity matrix entries is required.

% ---------------------------------------------------------------------------
\section{Conclusion}
% ---------------------------------------------------------------------------

We have presented AXIOM, a formally specified protocol for allocating capital to
autonomous agents via verifiable proof of value creation. The central
contributions are:

\begin{enumerate}[(1)]
  \item A formal definition of the Proof of Value Creation (PVC) as a
        risk-adjusted, capital-normalised, cryptographically attested record
        of net economic output.
  \item The AxiomScore allocation function, which monotonically combines
        performance, identity reputation, strategy reputation, and tail-risk
        penalty.
  \item Proof that the slashing mechanism renders fraud economically irrational
        under a mild stake condition (Theorem~\ref{thm:fraud}).
  \item Proof that the allocation converges to a fixed point under stationary
        agents (Theorem~\ref{thm:fixedpoint}) and that PVE remains positive
        under the fraud-irrationality regime (Theorem~\ref{thm:pve}).
  \item An attack analysis covering Sybil, fake PVC, reputation laundering,
        oracle manipulation, and correlation attacks.
\end{enumerate}

AXIOM and Proof Protocol form a natural two-layer stack: execution integrity
below, economic value allocation above. Together they constitute an
infrastructure for a world where autonomous agents are funded, evaluated, and
rewarded purely on the basis of verifiable, risk-adjusted value creation.

% ---------------------------------------------------------------------------
\bibliographystyle{plain}
\begin{thebibliography}{9}

\bibitem{cpp2025}
Proof Protocol Authors.
\emph{A Replay-Based Verifiable Execution Layer for Multi-Paradigm AI
Computation}.
Technical Report, 2025.
\texttt{https://github.com/haynbroit-alt/cpo/blob/main/paper/proof\_protocol.tex}

\bibitem{nakamoto2008}
S. Nakamoto.
\emph{Bitcoin: A Peer-to-Peer Electronic Cash System}.
2008.

\bibitem{wood2014}
G. Wood.
\emph{Ethereum: A Secure Decentralised Generalised Transaction Ledger}.
Ethereum Project Yellow Paper, 2014.

\bibitem{roughgarden2021}
T. Roughgarden.
\emph{Transaction Fee Mechanism Design for the Ethereum Blockchain}.
EC '21, 2021.

\end{thebibliography}

\end{document}
